\documentclass[11pt]{article}

\usepackage{comment}
\usepackage{algorithmicx}
\usepackage{amstext,amssymb,amsfonts,latexsym,amsmath,amscd,amsthm}
\usepackage{subfig}
\usepackage[export]{adjustbox}
\usepackage{graphicx}
\usepackage{multirow}
\usepackage[round]{natbib}
\usepackage{booktabs}
\usepackage{lscape}
\usepackage{verbatim}
\usepackage{array}
\usepackage{setspace}
\usepackage{appendix}
\usepackage{palatino}
\usepackage{algcompatible}
\usepackage{algorithm}
\usepackage{algpseudocode}
\usepackage{xurl}
\usepackage{color}
\usepackage{rotating}
\usepackage{longtable}
\usepackage{enumitem}
\usepackage{ulem}

\newtheorem{theorem}{Theorem}[section]
\newtheorem{lemma}[theorem]{Lemma}
\newtheorem{proposition}{Proposition}
\newtheorem{corollary}[theorem]{Corollary}
\newtheorem{observation}{Observation}
\newtheorem{definition}{Definition}

\algrenewcommand\algorithmicrequire{\textbf{Input:}}
\algrenewcommand\algorithmicensure{\textbf{Output:}}

\usepackage[
  pdfauthor={},
  pdftitle={Vocabulary Covering Optimization},
  pdfstartview=XYZ,
  bookmarks=true,
  colorlinks=true,
  linkcolor=blue,
  urlcolor=blue,
  citecolor=blue,
  pdftex,
  linktocpage=true,
  hyperindex=true
]{hyperref}

\usepackage{tikz}
\usetikzlibrary{arrows}
\usepackage{enumerate}
\usepackage{caption}
\bibliographystyle{apalike}

\textwidth = 6.5 in
\textheight = 9 in
\topmargin = -0.5 in
\evensidemargin = 0 in
\oddsidemargin = 0 in
\parskip = 0.1 in
\parindent = 0 in
\tabcolsep = 5 pt

\DeclareMathOperator{\rank}{rank}
\DeclareMathOperator*{\argmin}{arg\,min}
\DeclareMathOperator*{\argmax}{arg\,max}

% ─────────────────────────────────────────────────────────────────────────────
\title{\textbf{Frequency-Weighted p-Median\\
       on Semantic Word Space}\\[6pt]
       {\large EMÜ414 -- Mathematical Programming and Computer Applications\\
        Hacettepe University, Department of Industrial Engineering\\
        Spring 2025--2026}}

\author{[Team Name]\\[4pt]
        [Member 1] \quad [Member 2] \quad [Member 3]\\
        [Member 4] \quad [Member 5] \quad [Member 6]}

\date{May 2026}
% ─────────────────────────────────────────────────────────────────────────────

\begin{document}
\maketitle

% ─────────────────────────────────────────────────────────────────────────────
\section{Introduction}
% ─────────────────────────────────────────────────────────────────────────────

Vocabulary selection is a fundamental problem in language learning and natural language
processing. Given a large lexicon, a learner or system with limited capacity must choose
which words to focus on so that the chosen set represents the broader vocabulary as well
as possible. More frequent words are inherently more valuable to cover, since they appear
more often in real text.

In this project, we model vocabulary selection as a \textit{Frequency-Weighted p-Median Problem} (FWPMP) over a semantic word space. We take the 100 most
representative words drawn from the top 10{,}000 most frequent English words, build a
pairwise semantic distance matrix using pre-trained GloVe word embeddings, and select
a budget-constrained set of \emph{representative words} that minimises the total
frequency-weighted semantic distance between every word and its assigned representative.

The problem belongs to the class of \textbf{Binary Integer Programming (BIP)} problems.
It generalises the classical $p$-Median facility location problem to a weighted semantic
setting, and is known to be NP-hard \citep{hochbaum1997}.

% ─────────────────────────────────────────────────────────────────────────────
\section{Problem Statement}\label{sec:problem}
% ─────────────────────────────────────────────────────────────────────────────

\subsection{Informal Description}

We are given a set of $|W| = 100$ English words. Each word $w$ has a known usage
frequency weight $F_w$, derived from its rank in the top-10k English frequency list.
A pairwise semantic distance $D_{wv}$ is defined for every pair $(w, v)$ using the
cosine distance between GloVe embedding vectors.

We select at most $B$ words as \emph{representatives} and assign every remaining
word to exactly one representative. The goal is to minimise the total
frequency-weighted semantic distance between words and their assigned representatives
--- equivalently, to bring the most-used words as close as possible to a representative.

\subsection{Mathematical Formulation}\label{sec:model}

\subsubsection*{Notation}

\begin{table}[h]
\caption{Mathematical Notation}
\label{tab:notation}
\centering
\renewcommand{\arraystretch}{1.3}
\begin{tabular}{ll}
\toprule
\textbf{Symbol} & \textbf{Description} \\
\midrule
\multicolumn{2}{l}{\textit{Sets}} \\
$W$ & Set of words, $|W| = 100$, indexed by $w$ and $v$ \\[4pt]
\multicolumn{2}{l}{\textit{Parameters}} \\
$D_{wv} \in [0,2]$ & Semantic cosine distance between word $w$ and word $v$ \\
$F_{w} \in \mathbb{R}_{+}$ & Frequency weight of word $w$ (higher = more frequent) \\
$B \in \mathbb{Z}_{+}$ & Budget: maximum number of representatives to select \\[4pt]
\multicolumn{2}{l}{\textit{Decision Variables}} \\
$x_{v} \in \{0,1\}$ & $1$ if word $v$ is selected as a representative, $0$ otherwise \\
$y_{wv} \in \{0,1\}$ & $1$ if word $w$ is assigned to representative $v$, $0$ otherwise \\
\bottomrule
\end{tabular}
\end{table}

\subsubsection*{Integer Programming Model}

\begin{alignat}{3}
    &\underset{}{\text{minimize}}\; & &
    \sum_{w \in W}\sum_{v \in W} F_{w} \cdot D_{wv} \cdot y_{wv}
    \label{eq:obj} \\[6pt]
    &\text{subject to} \notag \\
    & & & \sum_{v \in W} y_{wv} \;=\; 1
    \qquad \forall\, w \in W
    \label{eq:assign} \\
    & & & y_{wv} \;\leq\; x_{v}
    \qquad \forall\, w, v \in W
    \label{eq:link} \\
    & & & \sum_{v \in W} x_{v} \;\leq\; B
    \label{eq:budget} \\
    & & & x_{v},\; y_{wv} \;\in\; \{0,1\}
    \qquad \forall\, w, v \in W
    \label{eq:binary}
\end{alignat}

\paragraph{Objective \eqref{eq:obj}.}
Minimise the total frequency-weighted semantic distance summed over all
word-to-representative assignments. A high-frequency word ($F_w$ large) assigned
far from its representative ($D_{wv}$ large) incurs a large cost, so the model
prioritises placing popular words close to a selected representative.

\paragraph{Assignment constraints \eqref{eq:assign}.}
Every word must be assigned to exactly one representative.
Since $D_{ww} = 0$, a selected word is always assigned to itself at zero cost.

\paragraph{Linking constraints \eqref{eq:link}.}
A word $w$ can only be assigned to word $v$ if $v$ has been selected as a
representative ($x_v = 1$).

\paragraph{Budget constraint \eqref{eq:budget}.}
At most $B$ words may be selected as representatives.

\paragraph{Binary constraints \eqref{eq:binary}.}
All decision variables are binary.

\subsubsection*{Problem Classification}

The model is a \textbf{Binary Integer Program (BIP)}.  It is an instance of the
weighted $p$-Median facility location problem, which is NP-hard in general
\citep{hochbaum1997}.
The distance matrix $D$ enters the objective directly, so the
actual semantic distances --- not merely a coverage threshold --- drive the solution.

% ─────────────────────────────────────────────────────────────────────────────
\section{Data}\label{sec:data}
% ─────────────────────────────────────────────────────────────────────────────

\subsection{Word Selection}

We start from the top 10{,}000 most frequent English words as ranked by the
\texttt{wordfreq} corpus \citep{speer2022}.
To ensure a diverse semantic spread, we divide the list into 10 frequency bands
of 1{,}000 words each and sample 10 words uniformly at random from each band
(random seed = 42). This stratified procedure yields a subset of 100 words that
spans both very common ($\text{rank} \leq 1000$) and moderately common
($\text{rank} \leq 10{,}000$) vocabulary.

\subsection{Frequency Weights}

Each word $w$ with frequency rank $\rho_w$ is assigned the weight
\begin{equation}
    F_w = \frac{1}{\rho_w},
    \label{eq:weight}
\end{equation}
so higher-frequency words have larger weights. Weights are normalised to
$[0,1]$ by dividing by $\max_w F_w = F_{\text{rank}=1}$.

\subsection{Semantic Distance Matrix}

Word embeddings are obtained from the pre-trained GloVe model with
100-dimensional vectors trained on the 6B-token Wikipedia and Gigaword corpus
\citep{pennington2014}.
The semantic distance between words $w$ and $v$ is defined as the cosine distance
\begin{equation}
    D_{wv} = 1 - \frac{\mathbf{e}_w \cdot \mathbf{e}_v}
                      {\|\mathbf{e}_w\|\,\|\mathbf{e}_v\|},
    \label{eq:dist}
\end{equation}
where $\mathbf{e}_w \in \mathbb{R}^{100}$ is the GloVe embedding of word $w$.
The resulting $100 \times 100$ distance matrix $D$ is symmetric with zeros on
the diagonal and values in $[0, 2]$.

% ─────────────────────────────────────────────────────────────────────────────
\section{Computational Results}\label{sec:results}
% ─────────────────────────────────────────────────────────────────────────────

[To be completed after running model.py]

Experiments are conducted with varying values of the budget $B \in \{5, 10, 15, 20\}$.
The model is solved to proven optimality using Gurobi 13.x on a standard laptop.
For each value of $B$, we report the optimal objective value (total frequency-weighted
distance) and the average distance per word assignment.

% ─────────────────────────────────────────────────────────────────────────────
\section{Sensitivity Analysis}\label{sec:sensitivity}
% ─────────────────────────────────────────────────────────────────────────────

[To be completed]

We analyse:
\begin{itemize}
    \item Effect of budget $B$ on total frequency-weighted distance (objective value).
    \item Diminishing returns: how much does each additional representative reduce cost?
    \item Assignment clusters: which words group together under the same representative,
          and does the clustering reflect semantic similarity?
\end{itemize}

% ─────────────────────────────────────────────────────────────────────────────
\section{Conclusion}\label{sec:conclusion}
% ─────────────────────────────────────────────────────────────────────────────

[To be completed]

% ─────────────────────────────────────────────────────────────────────────────
\bibliographystyle{apalike}
\begin{thebibliography}{9}

\bibitem{hochbaum1997}
Hochbaum, D.~S. (1997).
\newblock \emph{Approximation Algorithms for NP-Hard Problems}.
\newblock PWS Publishing.

\bibitem{nemhauser1978}
Nemhauser, G.~L., Wolsey, L.~A., and Fisher, M.~L. (1978).
\newblock An analysis of approximations for maximizing submodular set functions.
\newblock \emph{Mathematical Programming}, 14(1):265--294.

\bibitem{pennington2014}
Pennington, J., Socher, R., and Manning, C.~D. (2014).
\newblock Gloʋe: Global vectors for word representation.
\newblock In \emph{Proceedings of EMNLP}, pages 1532--1543.

\bibitem{speer2022}
Speer, R. (2022).
\newblock \texttt{wordfreq}: Access a database of word frequencies, in multiple languages.
\newblock \url{https://github.com/rspeer/wordfreq}.

\end{thebibliography}

\end{document}
